\documentclass[12pt,a4paper]{article}
\usepackage{fontspec}
\setmainfont{Liberation Serif}   % или Times New Roman, Arial и др. с кириллицей
\usepackage[russian]{babel}
\usepackage{amsmath,amssymb,amsfonts,mathtools}
\usepackage{geometry}
\geometry{left=2cm,right=2cm,top=2cm,bottom=2cm}
\usepackage{booktabs}
\usepackage{hyperref}
\usepackage{url}
\usepackage{graphicx}
\usepackage{caption}
\usepackage{subcaption}
\usepackage{float}
\usepackage{siunitx}
\usepackage{bm}
\usepackage{pgfplots}
\pgfplotsset{compat=1.18}
\usepackage{xcolor}
\usepackage{multicol}
\usepackage{tikz}
\usetikzlibrary{positioning, arrows.meta, calc}

\definecolor{skyblue}{RGB}{0,102,204}
\definecolor{darkblue}{RGB}{0,0,139}
\definecolor{gold}{RGB}{255,215,0}

% Макросы YUCT (оставлены без изменений)
\newcommand{\eps}{\varepsilon}
\newcommand{\kappac}{\kappa_c}
\newcommand{\Keff}{K_{\mathrm{eff}}}
\newcommand{\betaf}{\beta}
\newcommand{\df}{d_{\!f}}
\newcommand{\Mpl}{M_{\mathrm{Pl}}}
\newcommand{\Lpl}{l_{\mathrm{Pl}}}
\newcommand{\tP}{t_{\mathrm{P}}}
\newcommand{\Sodd}{S_{\mathrm{odd}}}
\newcommand{\Seven}{S_{\mathrm{even}}}
\newcommand{\Dtau}{\Delta\tau_{\min}}
\newcommand{\Lzero}{L_0}
\newcommand{\Omegarot}{\Omega_{\mathrm{rot}}}
\newcommand{\Mbar}{M_{\mathrm{bar}}}
\newcommand{\muniv}{M_{\mathrm{univ}}}

\title{\textbf{Критическая масса локального Большого взрыва как координационная основа YUCT: самодостаточное резюме}}
\author{Алексей В. Якушев}
\date{Апрель 2026}

\begin{document}
	
	\maketitle
	
	\begin{abstract}
		Настоящий документ представляет собой краткое, самодостаточное изложение ключевых результатов Объединяющей теории координации Якушева (YUCT), касающихся критической массы локального Большого взрыва, связанных инфляционных параметров, глобального вращения Вселенной и космической иерархии масс. Все уравнения явно помечены как соотношения YUCT. Документ структурирован для удобного восприятия как человеком, так и моделями ИИ: приведены чёткие формулы, таблицы и проверяемые предсказания. Представлена полная иерархия масс (включая решенную YUCT классическую проблему иерархии (17 порядков величины) от планковского масштаба до наблюдаемой Вселенной, демонстрирующая объединение квантовых, звёздных и космологических масштабов посредством единого координационного закона масштабирования.
	\end{abstract}
	
	\section{Фундаментальные постулаты YUCT}
	
	\begin{itemize}
		\item \textbf{Эффективность координации} (протокол YPSDC, Приложение B):
		\[
		\Keff = \frac{H(\mathcal{A})}{H(\kappa)}, \qquad\text{(YUCT-1)}
		\]
		где \(H(\mathcal{A})\) — энтропия активированного действия, \(H(\kappa)\) — энтропия переданного короткого индекса.
		
		\item \textbf{Универсальный закон масштабирования ошибки} (Приложение L):
		\[
		\varepsilon = \kappac \alpha \Keff^{-\betaf}, \qquad \betaf = \frac{2}{3},\ \kappac = \frac{1}{3}. \qquad\text{(YUCT-2)}
		\]
		
		\item \textbf{Алгебраическая петля} (Приложение Y):
		\[
		\Sodd = \frac{6}{5},\quad \Seven = \frac{4}{5},\quad \frac{\Sodd}{\Seven} = \frac{3}{2} = \frac{1}{\betaf}. \qquad\text{(YUCT-3)}
		\]
		Фундаментальные кванты длины и времени:
		\[
		\Lzero = \frac{2}{3}\Lpl,\qquad \Dtau = \frac{2}{3}\tP. \qquad\text{(YUCT-4)}
		\]
		
		\item \textbf{Фрактальная размерность} из информационного насыщения (Приложение L):
		\[
		d_f = \frac{11}{5}=2.2. \qquad\text{(YUCT-5)}
		\]
	\end{itemize}
	
	\section{Критическая масса локального Большого взрыва}
	
	Используя закон масштабирования, выведенный из непротиворечивости координационной сети,
	\[
	\Keff(R) = \left(\frac{R}{\Lzero}\right)^{\betaf d_f/2}, \qquad\text{(YUCT-6)}
	\]
	и критическое условие \(\varepsilon \approx 1\) (т.е. \(\Keff = K_{\mathrm{crit}} = 3^{3/2}\)), получаем критическую массу гравитационно связанной доколлапсной области:
	\[
	M_{\mathrm{crit}} = 3 M_{\mathrm{Pl}} \approx 6.5\times 10^{-8}\ \mathrm{kg}. \qquad\text{(YUCT-7)}
	\]
	Вывод использует шварцшильдовский радиус \(R_S = 2GM/c^2\) и алгебраическую петлю \(\Lzero = \frac{2}{3}\Lpl\). Квантовые поправки оставляют коэффициент 3 неизменным в главном порядке.
	
	\section{Космологическая инфляция как координационный фазовый переход}
	
	Инфляционные наблюдаемые параметры предсказываются из тех же постулатов:
	\begin{align}
		n_s &= 1 - 2\betaf^2 + \frac{4}{3}\betaf^3 + \cdots \approx 0.965, \qquad\text{(YUCT-8)}\\
		r &= 8(2\betaf)^2 = 32\betaf^2 \approx 0.07, \qquad\text{(YUCT-9)}\\
		f_{\mathrm{NL}}^{\mathrm{local}} &= \frac{5}{3}\betaf \approx 1.12. \qquad\text{(YUCT-10)}
	\end{align}
	Эти предсказания проверяемы с помощью экспериментов CMB‑S4, LiteBIRD, Euclid и SPHEREx.
	
	\section{Глобальное вращение Вселенной}
	
	Диполь реликтового излучения частично обусловлен глобальным жёстким вращением с угловой скоростью \(\omega\). Амплитуда диполя растёт с красным смещением \(z\) как:
	\[
	v_{\mathrm{dip}}(z) = v_{\mathrm{local}} + \omega\, r(z). \qquad\text{(YUCT-11)}
	\]
	Из роста, наблюдаемого в обзорах NVSS, Quaia и других, получаем
	\[
	\omega \approx 8.5\times 10^{-22}\ \mathrm{rad/s},\qquad
	T_{\mathrm{rot}} = \frac{2\pi}{\omega} \approx 2.3\times 10^{14}\ \mathrm{yr}. \qquad\text{(YUCT-12)}
	\]
	Безразмерный параметр вращения:
	\[
	\Omegarot \equiv \frac{\omega}{H_0} = \frac{\Seven}{\Sodd} \betaf^2 \mathcal{F}(d_f)\left(\frac{L_0}{R_H}\right)^{2\betaf} \approx 3.9\times 10^{-4}, \qquad\text{(YUCT-13)}
	\]
	а параметр Хаббла удовлетворяет \(H_0 = \betaf\,\omega\). Барионная масса Вселенной следует из вращения:
	\[
	\Mbar = \Omega_b\frac{\betaf^2 c^3}{G H_0} \approx (2.3\pm 0.2)\times 10^{51}\ \mathrm{kg}. \qquad\text{(YUCT-14)}
	\]
	
	\section{Полная иерархия масс от квантового до космического масштаба}
	
	В YUCT все массы от планковской до наблюдаемой Вселенной подчиняются единому координационному закону масштабирования. В таблице~\ref{tab:full-hierarchy} перечислены уровни иерархии, типичные массы, примеры и соответствующие соотношения YUCT. На рисунке~\ref{fig:hierarchy-scheme} схематически изображены дискретные ступени лестницы.
	
	\begin{table}[htbp]
		\centering
		\caption{Полная иерархия масс от квантового до космического масштаба (YUCT).}
		\label{tab:full-hierarchy}
		\scriptsize
		\begin{tabular}{l c c l}
			\toprule
			\textbf{Уровень} & \textbf{Типичная масса (кг)} & \textbf{Примеры} & \textbf{Соотношение масштабирования YUCT} \\
			\midrule
			Планковское семя & \(M_{\mathrm{crit}} = 3M_{\mathrm{Pl}} \approx 6.5\times10^{-8}\) & Семя локального Большого взрыва & \(M_{\mathrm{crit}} = 3 M_{\mathrm{Pl}}\) (YUCT-7) \\
			Электрон & \(m_e \approx 9.1\times10^{-31}\) & Легчайший заряженный фермион & \(m_e = M_{\mathrm{Pl}} (2/3)^{96+\sigma}\xi_e\) \\
			Протон & \(m_p \approx 1.67\times10^{-27}\) & Стабильный барион & \(m_p = M_{\mathrm{Pl}} (2/3)^{96+\sigma}\xi_p\) \\
			Белый карлик & \(\sim 10^{30}\) (\(0.6M_\odot\)) & Предел Чандрасекара & \(M_{\mathrm{WD}} \propto M_{\mathrm{Pl}}^{3}/m_p^2\) \\
			Нейтронная звезда & \(\sim 3\times10^{30}\) (\(1.4M_\odot\)) & Предел TOV & \(M_{\mathrm{NS}} \propto M_{\mathrm{Pl}}^{3}/m_p^2\) \\
			Сверхмассивная ЧД & \(\sim 10^{36}\)–\(10^{40}\) & Центры галактик & \(M_{\mathrm{BH}} \propto M_{\mathrm{Pl}} (2/3)^{N}\) \\
			Наблюдаемая Вселенная & \(M_{\mathrm{univ}} \approx 1.5\times10^{53}\) & Полная масса внутри \(R_H\) & \(M_{\mathrm{univ}} = \frac{\betaf^2 c^3}{G H_0}\) (YUCT-14) \\
			Барионная Вселенная & \(M_{\mathrm{bar}} \approx 2.3\times10^{51}\) & Видимое вещество & \(M_{\mathrm{bar}}/m_p = (M_{\mathrm{Pl}}/m_e)^{3.5}\) (YUCT-15) \\
			\bottomrule
		\end{tabular}
	\end{table}
	
	\begin{figure}[htbp]
		\centering
		\begin{tikzpicture}[
			level/.style={rectangle, draw=darkblue!70, fill=skyblue!10, thick, minimum width=7cm, minimum height=0.7cm, rounded corners, align=center, font=\small},
			arrow/.style={->, >=Stealth, thick, color=darkblue},
			]
			\node[level] (planck) {Планковское семя: \(M_{\mathrm{crit}} = 3M_{\mathrm{Pl}}\)};
			\node[level, below=0.6cm of planck] (electron) {Электрон: \(m_e\)};
			\node[level, below=0.6cm of electron] (proton) {Протон: \(m_p\)};
			\node[level, below=0.6cm of proton] (wd) {Белый карлик: \(\sim 0.6 M_\odot\)};
			\node[level, below=0.6cm of wd] (ns) {Нейтронная звезда: \(\sim 1.4 M_\odot\)};
			\node[level, below=0.6cm of ns] (bh) {Сверхмассивная ЧД: \(10^6\)–\(10^{10} M_\odot\)};
			\node[level, below=0.6cm of bh] (univ) {Наблюдаемая Вселенная: \(M_{\mathrm{univ}}\)};
			\draw[arrow] (planck) -- (electron);
			\draw[arrow] (electron) -- (proton);
			\draw[arrow] (proton) -- (wd);
			\draw[arrow] (wd) -- (ns);
			\draw[arrow] (ns) -- (bh);
			\draw[arrow] (bh) -- (univ);
		\end{tikzpicture}
		\caption{Дискретная лестница масс от планковского семени до Вселенной в YUCT. Каждая ступень соответствует резонансу координационного поля.}
		\label{fig:hierarchy-scheme}
	\end{figure}
	
	Космическая иерархия масс дополнительно выражается алгебраическим соотношением
	\[
	\boxed{\frac{\Mbar}{m_p} = \left(\frac{M_{\mathrm{Pl}}}{m_e}\right)^{\mathcal{S}}},\qquad
	\mathcal{S} \equiv \Sodd + \Seven + \frac{\Sodd}{\Seven} = 3.5. \qquad\text{(YUCT-15)}
	\]
	Численно \((M_{\mathrm{Pl}}/m_e)^{3.5} \approx 1.88\times 10^{78}\), что согласуется с наблюдаемым \(\Mbar/m_p \approx 1.4\times 10^{78}\) с точностью до множителя 1.3; расхождение объясняется эффективностью повторного нагрева \(\epsilon_{\mathrm{reh}}\approx 0.85\) и геометрическим фактором \(\mathcal{F}_{\mathrm{rot}}\approx 0.9\). Это демонстрирует, что те же алгебраические константы \(S_{\mathrm{odd}}, S_{\mathrm{even}}\), которые управляют универсальным показателем ошибки, определяют и барионный бюджет Вселенной.
	
	\section{Иерархия космических объектов и критическое вращение}
	
	Для самогравитирующих систем, поддерживаемых давлением вырождения, YUCT предсказывает универсальное масштабирование критической угловой скорости:
	\[
	\Omega_{\mathrm{crit}} \propto M^{\betaf},\qquad \betaf = 2/3. \qquad\text{(YUCT-16)}
	\]
	Таблица~\ref{tab:objects} и рисунок~\ref{fig:omega-mass} иллюстрируют эту иерархию от планет до наблюдаемой Вселенной.
	
	\begin{table}[htbp]
		\centering
		\caption{Критические параметры выбранных космических объектов (YUCT).}
		\label{tab:objects}
		\small
		\begin{tabular}{l c c c}
			\toprule
			\textbf{Объект} & \textbf{Масса (\(M_\odot\))} & \(\log_{10}(\Omega_{\mathrm{crit}}/\mathrm{s^{-1}})\) & \textbf{Закон масштабирования} \\
			\midrule
			Газовый гигант (Юпитер) & 0.001 & \(-3.8\) & \(R\sim M^{-0.67}\) \\
			Красный карлик (M5V) & 0.2 & \(-4.5\) & \(L\sim M^{2.3}\) \\
			Белый карлик & 0.6 & 0.2 & \(R\sim M^{-1/3}\) \\
			Нейтронная звезда & 1.4 & 3.7 & \(M\sim\Lambda^{0.67}\) \\
			Сверхмассивная НЗ & 2.2 & 4.2 & \(\Omega_{\mathrm{crit}}\sim M^{0.67}\) \\
			Звёздная чёрная дыра & 5 & -- & \(R_S\sim M\) \\
			Наблюдаемая Вселенная & \(7.5\times10^{22}\) & \(-17.7\) (\(H_0\)) & \(\rho_c\sim H^2\sim K_{\mathrm{eff}}^{2/3}\) \\
			\bottomrule
		\end{tabular}
	\end{table}
	
	\begin{figure}[htbp]
		\centering
		\begin{tikzpicture}
			\begin{loglogaxis}[
				xlabel={Масса \(M\) (\(M_\odot\))}, ylabel={\(\Omega_{\mathrm{crit}}\) (с\(^{-1}\))},
				grid=both, width=0.9\textwidth, height=0.6\textwidth,
				legend pos=north west, title={Критическое вращение компактных объектов (YUCT)},
				xmin=0.08, xmax=15, ymin=5e-2, ymax=2e4,
				]
				\addplot[domain=0.1:10, samples=2, thick, blue, dashed] {10^(0.67*log10(x)+0.5)};
				\addlegendentry{\(\Omega_{\mathrm{crit}}\propto M^{2/3}\) (БК)};
				\addplot[domain=1:3, samples=2, thick, green!60!black, dashed] {10^(0.67*log10(x)+3.7)};
				\addlegendentry{\(\Omega_{\mathrm{crit}}\propto M^{2/3}\) (НЗ)};
				\addplot[only marks, mark=*, mark size=3pt, red] coordinates {(0.6,1.6)};
				\addplot[only marks, mark=square*, mark size=3pt, green!60!black] coordinates {(1.4,5000) (2.2,15000)};
				\addplot[only marks, mark=triangle*, mark size=3pt, orange] coordinates {(0.2,3e-5)};
			\end{loglogaxis}
		\end{tikzpicture}
		\caption{Критическая угловая скорость в зависимости от массы (YUCT). Наклоны подтверждают \(\beta=2/3\).}
		\label{fig:omega-mass}
	\end{figure}
	
	\section{Проверяемые предсказания}
	
	\begin{enumerate}
		\item \textbf{Реликтовые первичные чёрные дыры:} \(f_{\mathrm{relic}} \approx 0.05\) от тёмной материи. Предсказание опровергается, если \(f_{\mathrm{relic}}>0.2\) или \(<0.01\) с достоверностью 95\%.
		\item \textbf{Инфляционные параметры:} \(n_s = 0.965\pm0.005\), \(r\approx 0.07\). Исключается, если \(r<0.01\) или \(n_s\) вне интервала \(0.960\pm0.010\).
		\item \textbf{Негауссовость:} \(f_{\mathrm{NL}}^{\mathrm{local}} \approx 1.12\). Опровергается, если измеренное значение \(<0.5\) или \(>1.8\) с уровнем значимости \(>3\sigma\).
		\item \textbf{Лабораторное отклонение от \(E=mc^2\):} \(\Delta E/E \sim \kappa_c\alpha K_{\mathrm{eff}}^{-\beta}\). Нулевой результат при чувствительности \(10^{-18}\) поставил бы YUCT под сомнение.
		\item \textbf{Период космического вращения:} \(T_{\mathrm{rot}} \approx 2.3\times10^{14}\) лет. Будущие обзоры (Euclid, SKA) уточнят \(\omega\).
	\end{enumerate}
	
	\section{Заключение}
	
	Постулаты YUCT объединяют критическую массу локального Большого взрыва, инфляционные наблюдаемые параметры, глобальное вращение Вселенной и полную иерархию масс от планковского семени до наблюдаемого космоса. Все предсказания количественны и проверяемы. Основные уравнения сведены выше, каждое отмечено как соотношение YUCT. Детальные выводы приведены в цитируемых приложениях.
	
	\section*{Доступность данных}
	
	Полные выводы и код: \url{https://github.com/Alexey-Yakushev-YUCT/YPSDC}
	
	\section*{Ссылки на приложения YUCT}
	\begin{itemize}
		\item Приложение L: Фрактальное координационное масштабирование ошибки – универсальный закон \(\beta=2/3\), \(d_f=11/5\).
		\item Приложение Y: Математические основы и вывод KPZ.
		\item Приложение P: Температурная зависимость гравитации.
		\item Приложение M: Инфляция как координационный фазовый переход.
		\item Приложение \(\Omega\): Глобальное вращение Вселенной.
		\item Приложение B: Протокол YPSDC.
	\end{itemize}
	
\end{document}